% ============================================================================
%   CHAPTER 3 — PROJECT CONTEXT AND SCOPE
%   Aivancity: 5-8 pages
%   Purely academic research framework (no company internship)
% ============================================================================

\chapter{Project Context and Scope}
\label{chap:context}

% ----------------------------------------------------------------------------
\section{Academic Research Framework}
\label{sec:ctx-framework}

This thesis was conducted within the Final Year Project (\emph{Projet de 
Fin d'\'Etudes}, PFE) of the PGE5 program at aivancity, from March to 
September 2026. Unlike most PFE projects at aivancity, which are 
structured around a six-month industry internship in a host company, the 
present project is a purely academic research contribution. This choice 
was made in agreement with the academic supervisor, Dr.\ Abdelhamid Mellouk, 
after a preliminary review of the literature identified a specific 
methodological gap --- the systematic auditing of large language model 
calibration in clinical trial matching --- that could be addressed with 
academic resources alone, without requiring hospital or industry partners.

The academic framing has three practical consequences for the study. 
First, the research question and the four pre-registered hypotheses 
were formulated to be testable with publicly available datasets and 
open-source models, avoiding any dependency on proprietary systems or 
restricted patient data. Second, the study operates on a purely 
retrospective basis: no prospective clinical validation is performed, 
and no recommendations for real patient enrollment are made. Third, the 
computational budget is limited to what a single graduate student can 
access on free-tier cloud platforms, which caps the model sizes at 
seven to nine billion parameters and shapes several of the design 
choices described in Chapter~\ref{chap:methodology}.

Regular supervision meetings with Dr.\ Abdelhamid Mellouk were held throughout the 
project. Progress was formalized through monthly reports covering the 
scoping phase (March), the state of the art review (April), the 
experimental pipeline design (May), the execution of the twelve model 
runs (June), the extended analyses on mode collapse and truly 
discriminative cases (July), and the writing phase (August--September).

% ----------------------------------------------------------------------------
\section{Origin and Evolution of the Research Question}
\label{sec:ctx-origin}

The initial research question, formulated during the scoping phase, 
focused on patient-reported uncertainty communication in clinical trial 
matching: how do LLMs express uncertainty to end users, and does this 
communication align with what patients expect? A month of preliminary 
literature review, however, revealed that this framing conflated two 
distinct research questions --- the human factors of uncertainty 
communication, and the underlying probabilistic reliability of the 
model outputs --- and that only the second could be rigorously 
addressed within the scope of a Master's thesis.

The question was therefore reformulated to focus exclusively on the 
probabilistic reliability of LLM outputs. This reformulation offered 
three advantages. First, it aligned the study with an established 
research tradition on classifier calibration, providing well-defined 
metrics (ECE, Brier score, reliability diagrams) and clear evaluation 
thresholds (Van Calster's $0.05$). Second, it made the hypotheses 
testable with pre-registered statistical procedures, protecting the 
analysis against confirmation bias. Third, it identified a specific 
methodological gap in the medical LLM literature: while several papers 
had evaluated the accuracy of medical LLMs on clinical benchmarks, no 
systematic audit of their calibration had been performed for the 
clinical trial matching task. This gap became the specific target of 
the study.

The final research question --- ``Are the eligibility scores produced 
by LLM-based patient--trial matching systems probabilistically reliable, 
and does RLHF post-training affect this reliability?'' --- is a 
narrower, more falsifiable version of the initial intent, and the 
tradeoff is a deliberate methodological choice.

% ----------------------------------------------------------------------------
\section{Dataset: TrialGPT-Criterion-Annotations}
\label{sec:ctx-dataset}

\subsection{Dataset Description}
\label{subsec:ctx-desc}

The empirical foundation of this study is the 
TrialGPT-Criterion-Annotations dataset, released by Jin 
et~al.~\cite{jin2024trialgpt} alongside their TrialGPT system in 
\emph{Nature Communications}. The dataset comprises 1{,}015 
patient--criterion pairs, where each pair associates a synthetic 
patient note with a single eligibility criterion drawn from a real 
clinical trial listed on ClinicalTrials.gov. Each pair is annotated 
independently by three physicians (whose consensus label serves as 
the ground truth) and by GPT-4 (used as an automated second opinion). 
Both annotation sources apply the same six-level ordinal scale, 
described in detail in Section~\ref{subsec:met-data-composition}.

The dataset is publicly available under a CC BY 4.0 license, hosted on 
Zenodo alongside the TrialGPT source code. Its release supports the 
reproducibility of the original TrialGPT evaluation and, by extension, 
the present audit.

\subsection{Justification of the Single-Dataset Choice}
\label{subsec:ctx-single}

An important scoping decision of this study is the deliberate reliance 
on a single benchmark for the \emph{primary} analysis. Three 
alternatives were considered and rejected as the primary corpus during 
the design phase. First, the TREC Clinical Trials 2021--2022~\cite{roberts2021trec, roberts2022trec} 
tracks provide relevance judgments for patient--trial pairs, but they 
operate at the trial level rather than the criterion level and use 
only three ordinal grades (definitely relevant, possibly relevant, 
irrelevant). Adapting the calibration protocol to this coarser 
granularity would preclude the fine-grained analysis of inclusion 
versus exclusion criteria that motivates several of our findings. 
Second, the SIGIR 2016 clinical trial task shares the same limitations 
as TREC. Third, hospital-derived cohorts (n2c2, MIMIC-derived 
subsets) would provide realistic patient data but require 
institutional partnerships and ethics approval that fall outside the 
academic scope of the project.

TrialGPT-Criterion-Annotations is, to our knowledge, the only publicly 
available benchmark that combines per-criterion granularity, a 
six-level ordinal scale, and an explicit inclusion/exclusion typology. 
These properties are jointly essential for the calibration analysis 
described in Chapter~\ref{chap:methodology}, which is why it is retained 
as the primary corpus. The coarser granularity that disqualifies TREC 
for the calibration hypotheses does not, however, prevent it from 
serving a distinct and valuable role: because the mechanistic analysis 
of mode collapse (Chapter~\ref{chap:results}) operates on the model's 
probability distributions rather than on the six-level scale, TREC can 
be used as an independent corpus to test whether the collapse behaviour 
and its mechanism generalize beyond the primary benchmark. We therefore 
use TREC Clinical Trials 2021--2022 as a secondary, external validation 
corpus, described in Section~\ref{subsec:met-data-external} and reported 
in Section~\ref{sec:res-external}.

% ----------------------------------------------------------------------------
\section{Ethical Considerations}
\label{sec:ctx-ethics}

The ethical framework of this study is straightforward due to the 
nature of the dataset. TrialGPT-Criterion-Annotations contains no 
real patient data: all clinical notes were synthesized by the 
original authors from publicly available templates on 
ClinicalTrials.gov, using clinical concept substitution to produce 
notes that are semantically realistic yet fully non-identifying. No 
personal health information is processed at any stage of the 
pipeline, and no institutional review board approval or ethics 
committee clearance was required.

The absence of real patient data does not eliminate all ethical 
considerations. The results reported in Chapter~\ref{chap:results} 
concern models that could plausibly be deployed in real clinical 
screening workflows, and the mode collapse behaviors documented in 
Section~\ref{sec:res-mode-collapse} would translate into systematic 
screening failures if left undetected in production. The methodology 
described in this thesis is therefore intended as a template for 
pre-deployment audits rather than as a certification of the models 
themselves for clinical use. Practitioners integrating LLMs into 
recruitment workflows retain full responsibility for validating the 
models on their own populations and for maintaining human oversight 
of eligibility decisions.

Beyond patient data, the study also complies with the licensing terms 
of the models under evaluation. All six models are released under 
permissive open-source licenses (Apache 2.0 for Mistral, custom 
research licenses for Llama-3.1 and Meditron, MIT for BioMistral) 
that permit non-commercial evaluation and redistribution of 
inference outputs. The prediction files archived in the project 
repository (Appendix~\ref{app:code}) fall within the scope of these 
licenses.

% ----------------------------------------------------------------------------
\section{Constraints and Resources}
\label{sec:ctx-constraints}

\subsection{Computational Constraints}
\label{subsec:ctx-compute}

All experiments were performed on the Kaggle free-tier compute 
environment, which provides access to a single Tesla~T4 GPU with 
16~GB of memory and a weekly quota of 30~GPU-hours. This constraint 
shaped two important design choices. First, the model sizes are 
capped at seven to nine billion parameters, which excludes larger 
medical LLMs such as Meditron-70B or Med42-v2-70B from the 
evaluation. Second, all models are quantized to 4-bit NF4 precision, 
which reduces the memory footprint from approximately 14~GB per 
model (in half-precision) to approximately 4~GB. The justification 
of these choices and their potential impact on the results are 
discussed in Sections~\ref{sec:met-models} and~\ref{sec:met-infrastructure}.

\subsection{Model Access Constraints}
\label{subsec:ctx-access}

Access to the six selected models is subject to licensing terms that 
occasionally require prior approval. Llama-3.1 requires acceptance of 
Meta's community license through the HuggingFace hub, which was 
obtained during the setup phase. Gemma and Med42-v2, considered but 
excluded from the final selection, are subject to more restrictive 
licenses that complicate reproducibility and were partly responsible 
for their exclusion (Section~\ref{subsec:met-excluded}). All models 
retained in the final evaluation are freely accessible for research 
use without additional approval.

\subsection{Time Constraints}
\label{subsec:ctx-time}

The PFE was conducted over six months, from March to September 2026. 
The scoping phase (March) established the research question and the 
literature review. The design phase (April--May) formalized the 
pre-registered hypotheses and the experimental matrix. The execution 
phase (June--July) ran the twelve model configurations and the 
downstream analyses. The writing phase (August--September) produced 
the present manuscript. The time constraint precludes prospective 
clinical validation and extension to larger model scales; external 
validation was limited to a single independent corpus (TREC Clinical 
Trials, Section~\ref{sec:res-external}) rather than the broader 
multi-dataset replication that would be desirable. These directions 
are discussed in Chapter~\ref{chap:conclusion}.

% ----------------------------------------------------------------------------
\section{Scope of This Work}
\label{sec:ctx-scope}

The scope of this thesis is delimited both by what it includes and by 
what it deliberately excludes. Within scope, the study contributes: a 
systematic calibration audit of twelve LLM configurations on a 
primary benchmark; a pre-registered test of the effect of RLHF 
post-training on calibration in the clinical domain; a comparison 
between generalist and medically specialized open-source models; a 
documentation of two novel failure modes (mode collapse and complete 
failure on discriminative cases) that had not been previously reported 
for these models, together with a mechanistic reinterpretation of mode 
collapse as a recoverable property of the decoding rule; an evaluation 
of post-hoc recalibration, re-thresholding, and conformal prediction 
as practical remedies; and an external replication of the collapse 
mechanism on the independent TREC Clinical Trials corpus.

Explicitly out of scope, and deferred to future work in 
Chapter~\ref{chap:conclusion}, are: broader external validation across 
multiple independent benchmarks beyond the single-corpus TREC 
replication reported here; prospective clinical validation involving 
real patient enrollment; extension to model sizes above nine billion 
parameters; extension to multi-agent debate architectures as a 
mitigation for mode collapse; fine-tuning or continued pretraining of 
the evaluated models; and multilingual evaluation on non-English 
clinical text.

This delimitation is intended to make the contribution of the thesis 
falsifiable and specific, rather than exhaustive. The methodology 
proposed here can be extended in each of the excluded directions, and 
several such extensions are outlined as concrete research programs in 
Chapter~\ref{chap:conclusion}.